problem:  
Let \(G\) be a step‑2 Carnot group with homogeneous dimension \(Q\), Carnot–Carathéodory distance \(d_{CC}\), and Haar measure \(\mathfrak{m}\). For \(p>1\), consider global weak solutions \(u:G\to\mathbb{R}\) of the horizontal \(p\)-Laplace equation  
\[
\operatorname{div}_H(|\nabla_H u|^{p-2}\nabla_H u) = 0,
\]
with \(|\nabla_H u|\in L^p_{\mathrm{loc}}(\mathfrak{m})\).  

**Question:**  
Does there exist a (possibly group-dependent) threshold exponent \(p^*(G)>Q\) such that the following *universal Liouville rigidity* holds: every global \(p\)-harmonic function \(u\) on \(G\) with at most polynomial growth, i.e. \( |u(x)|\le C(1+d_{CC}(x,e))^k \) for some finite \(k\), is necessarily a horizontal affine function (i.e. a degree‑1 homogeneous polynomial with respect to the stratified dilation) whenever \(p\ge p^*(G)\)?  

Equivalently, is there a transition of analytic rigidity above some nonlinear threshold \(p^*(G)\), depending only on the subelliptic structure of \(G\), beyond which no nontrivial \(p\)-harmonic functions of controlled growth exist except the affine ones?  

Why is it a "good" problem:  
This refined problem removes artificial quantities and introduces a mathematically intrinsic notion of a “Liouville threshold.” It targets the intersection between subelliptic PDE theory, the metric geometry of Carnot groups, and nonlinear potential theory. Determining whether such a threshold \(p^*(G)\) exists would uncover how the interplay between the homogeneous dimension \(Q\) and the nonlinearity \(p\) controls global rigidity of \(p\)-harmonic structures. A positive result would indicate a universal analytic rigidity principle in all step‑2 Carnot geometries, hinting that high‑order degenerate elliptic equations may share geometric invariants with the underlying group stratification. A negative result would reveal unexpectedly rich families of globally defined nonlinear potentials and hence deepen understanding of function theory on noncommutative metric spaces. The problem is conceptually simple—does nonlinear rigidity above a critical exponent occur?—but resolving it would almost certainly require new techniques connecting analysis, Lie group geometry, and geometric measure theory, making it a compelling, forward-looking research question.